\documentclass[10pt,letterpaper]{article}
\usepackage{amsmath}
\usepackage{amsfonts}
\usepackage{amssymb}
\usepackage[english]{babel}
\usepackage{breakurl}
\usepackage[superscript]{cite}
\usepackage{fancyhdr}
\usepackage{float}
\usepackage[margin=1in]{geometry}
\usepackage{graphicx}
\usepackage{hyperref}
\usepackage[utf8]{inputenc}
\usepackage{makeidx}
\usepackage{multicol}
\usepackage{nth}
\usepackage{longtable}

% Use Helvetica
\usepackage[scaled]{helvet}
\renewcommand\familydefault{\sfdefault} 
\usepackage[T1]{fontenc}

\usepackage{setspace}
\usepackage{siunitx}
\usepackage{svg}
\usepackage{subcaption}
\usepackage{tikz}
\usepackage{titling}
\usepackage{url}
\usepackage{kantlipsum}

\usepackage{parskip} % Paragraph skip - adds extra lineskip spacing
\setlength{\parskip}{0.7\baselineskip plus 2pt}

% Custom definitions
\newcommand{\doctitle}{Requirements Specification}
\newcommand{\docsubtitle}{Multirotor Computing Platform with FPGA Hardware Acceleration for Machine Learning}

% Custom commands
\newcommand{\ts}{\textsubscript}	% Subscript command %

% Use hyphans to break up urls
\def\UrlBreaks{\do\/\do-}

% PDF and href setup
% Hyper ref
\hypersetup{
	colorlinks=true,
	citecolor=black,
	linkcolor=black,
	filecolor=black,
	urlcolor=blue,
	pdftitle={\@title},
	bookmarks=true
}
\urlstyle{same}

% Page headings
\pagestyle{fancy}
\fancyhead[L]{\MakeUppercase{CPEN/ELEC 491}}
\fancyhead[R]{\textbf{Team 109}}
\fancyfoot{}
\fancyfoot[C]{\thepage}

% No paragraph indent
\parindent 0ex

% Meta
\author{
	Deutsch, Peter &
	\textit{me@peterdeutsch.ca}
	\\
	He, Muchen &
	\textit{i@muchen.ca}
	\\
	Hsueh, Arthur &
	\textit{ah11962@outlook.com}
	\\
	Wang, Meng &
	\textit{wzfftxwd@gmail.com}
	\\
	Wilson, Ardell &
	\textit{ardellw96@gmail.com}
}
\title{\doctitle}
\date{\today}

\makeatletter
\renewcommand{\maketitle}{
	\bgroup
	\setlength{\parindent}{0pt}
	\begin{flushleft}
		% Top spacing
		\vspace*{1in}

		% Title
		\textbf{\Huge{\@title}}\\
		\hrulefill

		% Subtitle
		\textbf{\huge{\docsubtitle}}
		
		\vspace*{0.5in}

		% Course number and team
		\textbf{\Large{CPEN/ELEC 491 Capstone Team 109}}\\
		\hspace*{0.1cm}
		\begin{tabular}[h]{|ll}
			\@author
		\end{tabular}

		\vspace*{0.25in}

		\textbf{Website}: https://capstone-skynet.github.io

		\vfill

		% Logo
		\hspace*{-0.3cm}\includegraphics[scale=0.5]{../assets/ece_logo.pdf}

		% Date
		\large{\@date}
	\end{flushleft}
	\egroup
}
\makeatother

% Begin Document
\begin{document}

% Title Page
\begin{titlepage}
	\maketitle
\end{titlepage}

% Revision history
%\thispagestyle{empty}
\section*{Revision History}
The full revision history and commited changes of the document can be found in the git repository history: \href{https://github.com/Capstone-Skynet/Capstone-Skynet.github.io}{https://github.com/Capstone-Skynet/Capstone-Skynet.github.io/commits/master}.

\begin{table}[H]
\begin{tabular}{*{4}{l}p{0.5\linewidth}}
\hline
Version \# & Initials & Release Date & Changeset & Changes Made \\ \hline

0.0 & PD & 2019-10-11 & \texttt{660e001} & Initial skeleton of the document.\\
0.1 & MH & 2019-10-11 & \texttt{6af9e8a} & Populate initial document with draft content required for Milestone I.\\
0.2 & PD & 2019-11-23 & & Initial framework for test descriptions created.\\
0.3 & MW & 2019-11-24 & & First set of tests added.\\
1.0 & PD & 2019-11-25 & & General clean-up and release for Milestone II.\\
1.1 & AH & 2020-02-09 & & Added Full System Tests \\

 & & & \\ \hline
\end{tabular}
\end{table}

% Table of contents
\setcounter{secnumdepth}{3}
\tableofcontents

% Terms and Abbreviations
\thispagestyle{empty}

\section*{Terms and Abbreviations}

Technical terms and abbreviations dictionary go here.

\begin{tabular}[h]{rp{0.75\linewidth}}
    \hline
    \textbf{Term} & \textbf{Definition}\\
    \hline

    ANN & Artificial Neural Network, or simply ``Neural Network'', is a data processing model modeled after neuron interactions. The process consists of forward propagation using several matrix multiplications.\cite{ann}\\
    ASIC & Application-specific Integrated Circuits.\\
    CNN & Convolutional neural network are neural networks that is especially useful for image classification.\cite{cnn} \\
    ECE & Department of Electrical and Computer Engineering at the University of British Columbia.\\
    FPGA & Field-Programmable-Gate-Arrays, ``programmable'' hardware that allows ASIC-like performance with software-like turn-around time and flexibility.\\
    GPU & Graphics Processing Unit, a discrete piece of hardware designed to accelerate graphic-intensive or other parallel computing tasks.\\
    LOS & Line-of-sight.\\
    ML & Machine learning.\\
    Multirotor & An unmanned vehicle with multiple engines. \\
    OTS & Off-the-shelf, or commercially available/purchasable \\
    PID / PID Controller & Proportion-Integral-Derivative controllers is the most common control algorithm for precise and accurate movement, as well as to compensate external forces.\cite{pid}\\
    RNN & Recurrent neural networks are neural networks where the output depends on previous computations, essentially consists of memory.\cite{rnn}\\
    RX & Receiver.\\
    TC & Transport Canada.\\
    TX & Transmitter.\\
    YOLO & You-Only-Look-Once is a fast ML algorithm that detect objects but is unlike CNN nor RNN.\cite{yolo}\cite{yolo-2}\\
     & \\

    \hline

\end{tabular}


% List of figures and tables
\thispagestyle{empty}
\listoffigures
\listoftables
\newpage

% Set page and section counter
\setcounter{page}{1}

% Begin content
\section{About This Document}\label{section:about}
\subsection{Purpose}\label{section:about:purpose}
\subsection{Intended Audience}\label{section:about:audience}
\section{Context}\label{section:context}
\section{Domain}\label{section:domain}

\section{Goal}\label{section:goal}
The ultimate goal of this project is to develop a mobile computing platform which utilizes an FPGA to perform hardware-accelerated machine learning tasks and deploy it on an unmanned multirotor aerial system (drone). 

The drone should be piloted manually within the line-of-sight (LOS) of the pilot, using a ground-to-air transmitter (TX) in the form of an off-the-shelf radio controller and receiver combination. The drone should feature a flight controller capable of self-stabilization using well-tuned PIDs and,
with the payload attached, the drone's flight duration should be at least 50\% of that without the payload (typically 10--15 minutes).
The total takeoff mass of the integrated system should not exceed 25 kilograms --- as specified by Transport Canada, as a pilot with a \textit{Basic} or \textit{Advanced Operations} certificate cannot operate a drone heavier than 25kg.

The machine learning model should analyze the video stream from the onboard camera and detect pedestrians within the frame with near-humanlike accuracy. The model should exploit the FPGA's inherent parallelism to accelerate key ML tasks such as matrix multiplication and convolution. The model should output screen-space location of the detected pedestrians, with this prediction having a tolerable update latency of one second or less. In addition, the base station should display \textit{bounding boxes} surrounding the model's pedestrians predictions, overlaying a video of reasonable quality (at least 640x480, 10fps).


\section{Functional Specifications}\label{section:func-spec}
The main objectives are sorted by the client's priority in descending order.

\textbf{Integration of FPGA/Computer Vision Hardware with a Drone}:
As the ultimate deliverable consists of a highly mobile computing platform,
the student team's utmost objective is the successful integration of computation hardware with the drone.
The equipment for capturing video data, performing processing, and transmitting digital data must be reasonably compact and easily deployable.

\textbf{Air-to-Ground Data Transmission}:
Critical to research and analytics, the drone payload must be able to transmit video sensor data, in addition to processed data
from the onboard hardware accelerator (FPGA). There must be a ground station, consisting of a receiver and a display device (such as a mobile phone or 
computer) which receives and decodes the transmission for display, providing the user with a real-time video stream overlaid with machine learning intents.


\textbf{Machine Learning Implementation on the FPGA}:
The hardware-accelerated computing platform must correctly interface with attached peripherals (such as the camera sensor, the data transmitter (DTX), and other necessary hardware) through the use of well-defined communication protocols, such as serial or PCI-E (whichever is suitable given bandwidth requirements). The hardware accelerator must feature a starter ML model, implementing existing ML structures/frameworks such as (but not limited to) CNNs, RNNs, or YOLO \cite{yolo}, being fully capable of basic computer vision tasks.

\section{Non-Functional Specifications}\label{section:non-func-spec}

\section{Constraints}\label{section:const}
The most pressing constraints are of an administrative nature. The project described
herein is very large in scope --- the constrained 8-month development period 
will likely affect the team's maneuverability and ability to comprehensively mitigate risks. 
In addition, navigating regulatory hurdles regarding drone piloting and radio data transmission (set out by 
Transport Canada and Industry Canada, respectively) will incur additional time and resource constraints.

As this device is intended to be used as a development platform in the future, it is important that it is future-proofed for further modifications. As detailed in Section \ref{section:businessbackground}, FPGAs are
inherently constrained with regards to on-chip resources. An appropriately-sized FPGA will need to be selected, capable of
implementing the base-level ``glue-logic'' required for the platform to operate (for example, the camera interface and data transmission logic), while still
leaving enough available resources on the FPGA for the client to implement their own ML models after the project's delivery. 
This constraint is tightly coupled with the high level of legacy risk inherent in this project, as outlined in Section \ref{risk}.

The physical properties of the payload (FPGA, electronics, camera, and  DTX) is an additional constraint. Since the purpose of the project
is to validate the viability of mobile ML designs, the payload should be light, compact, and power efficient, such
that it can be carried on a drone for a reasonable flight time in order to carry out useful ML applications.
The weight of the payload (and its required batteries) affects total flight duration, maximum altitude, airspeed, and maneuverability. 

As the bandwidth of video transmission is limited by transmission frequency and the transmitter (DTX) itself, 
critical communication-related design decisions must be made early in the project.
A transmitter with slightly longer range would significantly increase the power consumption due to inverse-square law.
Moreover, the data throughput depends heavily on a successful implementation of the communication
protocols the student team decides to choose. Clever compression or encoding techniques might need to be
employed to enhance the throughput and achieve a reasonable result.



% Bibliography
\clearpage
\addcontentsline{toc}{section}{References}
\bibliographystyle{ieeetr}
\bibliography{references}

% Appendix (uncomment to enable appendix)
% \clearpage
% \appendix
% \section{Appendix name}\label{appendix:sample-appendix}
% Content here

\end{document}